% BEGIN THERMAL_R03_SYNC_18
\section{Evaluación de las constantes y contraste cuantitativo}
\label{sec:ii-contraste-cuantitativo}

Las construcciones anteriores permiten separar tres operaciones:
producir una magnitud, expresar su coordenada en una coordenatización y contrastar esa
coordenada con una referencia. El estado APP--TRIT--TPK conserva hojas,
orientación, cociente, acarreo y memoria antes de estas lecturas. Su
prolongación nonádica determina las coordenadas correlacionadas
\(\pi,\varphi,e,\alpha\); la norma multiplicativa sobre el cociente local y el retorno regional
producen las secciones de acción; la incidencia marcada \(90/120\) y
los canales conjugados producen el funcional de Barbero; el transporte
angular determina después las coordenadas CKM. Ninguna entrada de las tablas
externas interviene en la elección de esas operaciones.

Aquí se reúnen las evaluaciones de las construcciones anteriores y sus
referencias de comparación. Los dígitos de las salidas internas expresan
precisión de cálculo, no una incertidumbre experimental. Para un valor
\(x\) y una referencia no nula \(x_0\), ambos en la misma coordenatización, se usa
\[
 \Delta x=x-x_0,\qquad r_x=\frac{x-x_0}{x_0}.
\]
Cuando la referencia posee incertidumbres marginales asimétricas
\(u_-,u_+\), definimos además el residuo descriptivo
\begin{equation}
 z_x=\begin{cases}
 \Delta x/u_+,&\Delta x\geq0,\\
 \Delta x/u_-,&\Delta x<0.
 \end{cases}
 \label{eq:ii-residuo-marginal}
\end{equation}
Esta normalización no presupone independencia entre filas ni constituye
por sí sola una prueba estadística conjunta.

\subsection{Constante reducida de Planck y acción por vuelta: dos orientaciones}

En la coordenatización \(\mathcal U_S\mapsto\mathrm{J\,s}\), las secciones ya
construidas son
\begin{align}
 \hbar_{\rm pre}&=H_5\,10^{-34}\,\mathrm{J\,s},&
 H_5&=1.0545718176461544696\ldots,\\
 \hbar_{\rm ret}&=\eta_{\rm ret}\,10^{-34}\,\mathrm{J\,s},&
 \eta_{\rm ret}&=1.0545718169150390611\ldots,\\
 h_a&=2\pi\hbar_a,&a&\in\{\mathrm{pre},\mathrm{ret}\}.
 \label{eq:ii-cuatro-coordenadas-planck}
\end{align}
La década procede de \(\Sigma_H=\varphi\alpha^{16}\); la diferencia entre
las dos secciones procede del retorno regional. Multiplicar por \(2\pi\)
cambia de acción por radián a acción por vuelta, sin introducir otra
orientación ni alterar el residuo relativo.

La referencia del SI es
\[
 h_{\rm SI}=6.62607015\times10^{-34}\,\mathrm{J\,s},\qquad
 \hbar_{\rm SI}=\frac{h_{\rm SI}}{2\pi}
 =1.0545718176461563912\ldots\times10^{-34}\,\mathrm{J\,s}.
\]
Ambas expresiones son exactas. La expansión infinita de \(\hbar_{\rm SI}\)
no es una incertidumbre de medida. La fijación del valor numérico de
\(h\) pertenece a la definición del SI vigente desde 2019
\cite{bipm2018}; el contraste se realiza después de la construcción HMT.

\begin{table}[htbp]
\centering\small
\setlength{\tabcolsep}{5pt}
\begin{tabular}{@{}llr@{}}
\toprule
Magnitud & Coordenada en \(10^{-34}\,\mathrm{J\,s}\) & \(r_x\) frente al SI\\
\midrule
\(\hbar_{\rm pre}\) & \(1.0545718176461544696\) & \(-1.822221\times10^{-15}\)\\
\(\hbar_{\rm ret}\) & \(1.0545718169150390611\) & \(-6.932836\times10^{-10}\)\\
\(h_{\rm pre}\) & \(6.62607014999998793\) & \(-1.822221\times10^{-15}\)\\
\(h_{\rm ret}\) & \(6.62607014540625433\) & \(-6.932836\times10^{-10}\)\\
\bottomrule
\end{tabular}
\caption{Acción anterior y posterior al retorno. Las últimas cifras son
evaluaciones de las coordenadas obtenidas; los residuos no son errores
experimentales de \(h\) ni de \(\hbar\).}
\label{tab:ii-planck-comparacion}
\end{table}

\begin{proposition}[Invariancia del contraste bajo la vuelta completa]
En una coordenatización común,
\[
 \frac{h_a-h_{\rm SI}}{h_{\rm SI}}
 =\frac{\hbar_a-\hbar_{\rm SI}}{\hbar_{\rm SI}}.
\]
\end{proposition}
\begin{proof}
Las dos diferencias del primer miembro contienen el factor común
\(2\pi\), que se cancela con el del denominador.
\end{proof}
La separación pre/ret es cuantitativamente relevante: sustituir una por
otra cambia la comparación en varios órdenes de magnitud. El corrector
regional debe permanecer visible y no confundirse con redondeo numérico.

\subsection{Barbero--Immirzi: valor interno y comparación entre construcciones}

La evaluación del funcional se reconstruye con sus tres contribuciones:
\begin{align*}
 \Delta S_{90}&=2.95169439297457621139847987861\times10^{-4},\\
 \Delta S_{120}&=1.96835112502951720093738706217\times10^{-5},\\
 \pi AC^*/180&=0.270485322934140078465586458790\ldots,\\
 \gamma_{\rm HMT}&=\pi AC^*/180+12\Delta S_{90}-\Delta S_{120}\\
 &=0.27400767269445927474725526077398041940\ldots.
\end{align*}
Esta composición conserva el peso orientado \(12:-1\), anterior a su
evaluación decimal. Los sumandos indicados reconstruyen la coordenada
final con un residuo absoluto menor que \(3\times10^{-31}\).

Como comparación bibliográfica se considera el conteo de Ghosh y Mitra,
no una medición del parámetro. Su ecuación (7), escrita con
\(\lambda=2\pi\gamma_{\rm GM}\), es
\begin{equation}
 \sum_{j\in\{1/2,1,3/2,\ldots\}}(2j+1)
 e^{-2\pi\gamma_{\rm GM}\sqrt{j(j+1)}}=1.
 \label{eq:ii-gm-comparador}
\end{equation}
La referencia es A.~Ghosh y P.~Mitra,
\emph{An improved estimate of black hole entropy in the quantum geometry
approach}, \emph{Physics Letters B} \textbf{616} (2005), 114--117,
\href{https://arxiv.org/abs/gr-qc/0411035}{arXiv:gr-qc/0411035},
ecuaciones (7)--(8). La raíz se evalúa aquí para esa ecuación concreta;
no se identifica el conteo de puncturas con el generador HMT.

\begin{table}[htbp]
\centering\small
\begin{tabular}{@{}lll@{}}
\toprule
Construcción & Valor adimensional & Naturaleza\\
\midrule
Funcional \(\Gamma_{6\to8}\) & \(0.2740076726944593\) & Incidencia y retorno HMT\\
Raíz de \eqref{eq:ii-gm-comparador} & \(0.2740668582328300\) & Conteo bibliográfico declarado\\
\midrule
\(\gamma_{\rm HMT}-\gamma_{\rm GM}\) & \(-5.918553837\times10^{-5}\) & Diferencia entre construcciones\\
\(\gamma_{\rm HMT}/\gamma_{\rm GM}-1\) & \(-2.159529202\times10^{-4}\) & Diferencia relativa\\
\bottomrule
\end{tabular}
\caption{Contraste de Barbero--Immirzi en la orientación positiva. No se
asigna incertidumbre experimental ni significación estadística a este cuadro.}
\label{tab:ii-barbero-comparacion}
\end{table}

La evaluación de la raíz externa es reproducible sin truncar
arbitrariamente su cola. Con \(n=2j\), \(N=128\) y
\(r=e^{-0.27\pi}\), para \(\gamma\geq0.27\),
\[
 \sum_{n>N}(n+1)e^{-\pi\gamma\sqrt{n(n+2)}}
 \leq \frac{r^{N+1}[(N+2)-(N+1)r]}{(1-r)^2}<10^{-43}.
\]
Cada sumando disminuye estrictamente con \(\gamma\). Los extremos
\(0.27\) y \(0.28\) encierran la raíz; la bisección y la cota de cola
controlan los decimales de la tabla. Esta es una comprobación del
comparador bibliográfico, no una vía para seleccionar \(\gamma_{\rm HMT}\).

\subsection{Ángulos CKM, amplitudes y violación CP}

El lector angular ya construido determina, en grados,
\begin{equation}
 (\theta_{12},\theta_{23},\theta_{13},\delta)
 =(13.003313589509,\;2.398056157332,\;
 0.213977382244,\;65.676173123554)^\circ.
 \label{eq:ii-ckm-decimales}
\end{equation}
Estos valores proceden de \(M_{\rm CKM}(A,C^*/6,\Delta)^T\) y
\(\delta=9A\). La conversión a radianes se realiza una sola vez antes
de evaluar senos y cosenos. La parametrización unitaria de la sección
angular da
\[
 |V_{\rm CKM}|\simeq
 \begin{pmatrix}
 0.974350259&0.225005836&0.003734601\\
 0.224873110&0.973489279&0.041841465\\
 0.008582400&0.041121745&0.999117283
 \end{pmatrix},
\]
y el cuarteto invariante conserva la fase:
\[
 J=c_{12}c_{23}c_{13}^{\,2}s_{12}s_{23}s_{13}\sin\delta
 =3.1189723326632974\times10^{-5}\quad
 \text{en la evaluación del lector angular.}
\]
Esta evaluación utiliza \eqref{eq:ii-jarlskog} con las coordenadas de
\eqref{eq:ii-ckm-decimales}. Los nueve módulos impresos están redondeados; la unitariedad pertenece
a la matriz compleja anterior al redondeo, no a una matriz formada
solamente por módulos.

El comparador es el ajuste de tres generaciones del
\href{https://pdg.lbl.gov/2025/reviews/rpp2025-rev-ckm-matrix.pdf}
{Particle Data Group, actualización de 2025, \emph{CKM Quark-Mixing Matrix}},
ecuaciones (12.27)--(12.28) y el valor de \(J\) contiguo. Para evitar
alterar sus incertidumbres por una transformación no lineal, la tabla
conserva los observables expresados por PDG: \(s_{ij}=\sin\theta_{ij}\),
\(\delta\) en radianes y \(J\).

\begin{table}[htbp]
\centering\small
\setlength{\tabcolsep}{4pt}
\begin{tabular}{@{}lrrr@{}}
\toprule
Observable & Evaluación HMT & PDG 2025 & \(z_x\)\\
\midrule
\(s_{12}\) & \(0.225007404757\) & \(0.22501\pm0.00068\) & \(-0.003817\)\\
\(s_{23}\) & \(0.041841757010\) & \(0.04183^{+0.00079}_{-0.00069}\) & \(0.014882\)\\
\(s_{13}\) & \(0.003734601164\) & \(0.003732^{+0.000090}_{-0.000085}\) & \(0.028902\)\\
\(\delta\) (rad) & \(1.146265461116\) & \(1.147\pm0.026\) & \(-0.028251\)\\
\(10^5J\) & \(3.118972333\) & \(3.12^{+0.13}_{-0.12}\) & \(-0.008564\)\\
\bottomrule
\end{tabular}
\caption{Contraste marginal en la parametrización y edición declaradas.
El factor \(10^5\) multiplica conjuntamente valor e incertidumbre de la
última fila; no modifica su residuo normalizado.}
\label{tab:ii-ckm-comparacion}
\end{table}

Los residuos se calculan con las coordenadas anteriores al redondeo
de la tabla. El ajuste PDG combina datos y restricciones de unitariedad;
sus filas no son mediciones independientes. Por ello no se suma
\(\sum z_x^2\) ni se asigna un número de grados de libertad a este cuadro.
El resultado es una comparación marginal reproducible con esa edición,
no una medida de probabilidad del modelo. Las desviaciones marginales
históricas de los cálculos antecedentes no se trasladan a esta tabla:
su documentación no identifica la edición y la covarianza de aquel contraste.

El sector PMNS mantiene el transporte \(U=F_\ell^\dagger F_\nu\), los
dominios espectrales y los datos de fase desarrollados en la sección
angular. Una fila numérica necesita especificar los marcos del registro
empleado y la realización comparada. Esta tabla no atribuye cifras CKM
al lector leptónico ni determina una nueva matriz PMNS.

\subsection{Boltzmann y gravitación universal: magnitud, coordenatización y referencia}

La construcción térmica ya recuperada determina el transductor positivo
\(k_B:\mathcal L_\Theta\to\mathcal L_E\) y la relación
\begin{equation}
 k_B\Theta_{{\rm clk},a}\log3
 =\frac{h_a}{108t_0}=\frac{\hbar_a}{t_P},
 \qquad t_P=\frac{54}{\pi}t_0.
 \label{eq:ii-kb-comparacion}
\end{equation}
La información local, la acción y el período aparecen antes de la coordenatización
kelvin--joule. Para bases positivas \(u_E,u_\Theta\), si
\(k_B(u_\Theta)=b\,u_E\),
\(\Theta_{{\rm clk},a}=\theta_a u_\Theta\) y
\(E_{P,a}=\epsilon_a u_E\), la ecuación equivale a
\(b\theta_a=\epsilon_a/\log3\). Un cambio
\(u'_E=s_Eu_E\), \(u'_\Theta=s_\Theta u_\Theta\) produce
\[
 b'=\frac{s_\Theta}{s_E}b,\qquad
 \theta'_a=\frac{\theta_a}{s_\Theta},\qquad
 \epsilon'_a=\frac{\epsilon_a}{s_E}.
\]
La realización marcada de \ref{thm:ii-ter27-transducer} determina
además el coeficiente mediante la referencia espectral y la pareja
información por referencia--energía condicionada:
\begin{equation}
 b_*=10^{-23}\left(1+380\,10^{-3}+649\,10^{-6}\right)
     =\frac{1380649}{10^{29}}=1.380649\times10^{-23}.
 \label{eq:ii-b-evaluacion-comparada}
\end{equation}
Los grados, sus multiplicidades y el origen de capacidad se conservan
en \ref{sec:ii-ter27-antecedentes}; la caracterización aditiva y el
instrumento conservativo se desarrollan en
\ref{sec:ii-aditiva-instrumento}. Con $H$ y la referencia fijos,
$d\mathcal S/d\mathcal E=b_*\beta$ produce
$\Theta_P=(b_*\beta)^{-1}$ y $B_*(u_\Theta)=b_*u_E$.
Por tanto, la tabla muestra el coeficiente efectivamente evaluado,
junto con la referencia SI
\[
 k_B^{\rm SI}=1.380649\times10^{-23}\,\mathrm{J\,K^{-1}},
\]
El numeral coincide exactamente en esta realización y en las bases
transportadas indicadas. Su prueba termométrica comprende tanto la medida
espectral como los dos funcionales: cambiar sólo su normalización cambia
el lector. La identificación de $u_E/u_\Theta$ con $\mathrm{J/K}$ es
el transporte dimensional que permite la comparación. El valor SI es
definitorio; por ello no se le asigna incertidumbre experimental.

En gravedad, la realización circular fija
\begin{equation}
 G_a=\vartheta_{108}^{\circ}\frac{c_{\rm int}^{5}t_0^2}{\hbar_a},
 \qquad
 \vartheta_{108}^{\circ}=(54/\pi)^2
 =295.4525715010569\ldots.
 \label{eq:ii-g-comparacion}
\end{equation}
El selector y la reciprocidad \(G_a\hbar_a=
\vartheta_{108}^{\circ}c_{\rm int}^{5}t_0^2\) son cantidades internas de
la realización declarada. Una evaluación dimensional de \(G_a\) debe
conservar, además de la orientación \(a\), las coordenatizaciones de tiempo,
velocidad y acción. La composición longitudinal y radial de la
sección~\ref{g26:sec:rigidez-radial-es} determina primero
\[
 L=\frac{5759}{23040}\alpha_{\HMT}^{16}L_*,\qquad
 \hbar_{\rm ret}=\eta_{\rm ret}S_*,\qquad
 t_0=\frac{\pi_{\HMT}L}{54c_{\rm int}}.
\]
Con la misma sección de velocidad constitutiva, la sustitución en
\eqref{eq:ii-g-comparacion} da
\[
 G_{\rm ret}=\frac{c_{\rm int}^3L_*^2}{S_*}
 \left(\frac{5759}{23040}\right)^2
 \frac{\alpha_{\HMT}^{32}}{\eta_{\rm ret}}.
\]
La velocidad dimensional conserva la normalización constitutiva de los
canales angulares. Con las coordenadas ya generadas \(A_{\rm deg}\)
y \(C^*_{\rm deg}\), se escribe
\begin{equation}
 \begin{aligned}
 q_\pm&=\exp\!\left[-\frac{\pi_{\HMT}}{180}
               (A_{\rm deg}\pm C^*_{\rm deg})\right],\\
 r_\pm&=\frac{1-q_\pm^{90}}{1-q_\pm^{120}},\qquad
 \widehat c=\frac1{r_+r_-},\qquad c_{\rm int}=c_*\widehat c.
 \end{aligned}
 \label{eq:ii-g26-velocidad-constitutiva}
\end{equation}
En bases positivas compatibles,
\(\varepsilon=\varepsilon_*r_+^2\), \(\mu=\mu_*r_-^2\) y
\(c_*=(\mu_*\varepsilon_*)^{-1/2}\); por multiplicación,
\(c_{\rm int}=(\mu\varepsilon)^{-1/2}\). Así,
\begin{equation}
 G_{\rm ret}=\frac{c_*^3L_*^2}{S_*}
 r_N^2\alpha_{\HMT}^{32}\frac{\widehat c^{\,3}}{\eta_{\rm ret}}
 =\frac{L^2}{\hbar_{\rm ret}(\mu\varepsilon)^{3/2}}.
 \label{eq:ii-g26-gravedad-constitutiva}
\end{equation}
El evaluador utiliza la sección dimensional \(c_{\rm int}\) que
se declara a continuación y restituye \(c_*=c_{\rm int}/\widehat c\).
La sustitución verifica la equivalencia de estas expresiones en la misma
representación: \(\widehat c\) interviene una sola vez al pasar de
\(c_*\) a \(c_{\rm int}\) \cite{hmtg26counter}.

En la representación declarada \(L_*=1\,\mathrm m\),
\(S_*=10^{-34}\,\mathrm{J\,s}\) y
\(c_{\rm int}=299792458\,\mathrm{m\,s^{-1}}\), se obtiene
\[
 L=1.616254982625126330168262501334562\ldots\times10^{-35}
 \,\mathrm m,
\]
\[
 G_{\rm ret}=6.674299659414022706367776189\ldots\times10^{-11}
 \,\mathrm{m^3\,kg^{-1}\,s^{-2}}.
\]
El valor de comparación \(6.67430(15)\times10^{-11}\) se conserva
separadamente. La diferencia es
\(-3.40585977\ldots\times10^{-18}\,\mathrm{m^3\,kg^{-1}\,s^{-2}}\),
aproximadamente \(-0.00227057\) incertidumbres estándar de esa referencia.
Las cifras estables del cálculo describen la evaluación de la composición
con sus lectores y bases explícitos; la incertidumbre experimental pertenece
a la referencia medida.

\begin{table}[htbp]
\centering\small
\begin{tabularx}{\linewidth}{@{}lXX@{}}
\toprule
Objeto & Evaluación interna & Referencia dimensional\\
\midrule
\(B_*\) & \(1.380649\times10^{-23}\,u_E/u_\Theta\),
por \eqref{eq:ii-b-evaluacion-comparada} y la regla termométrica marcada.
& \(1.380649\times10^{-23}\,\mathrm{J/K}\), exacto en el SI.\\[4pt]
\(G_{\rm ret}\) & \(6.6742996594140227\ldots\times10^{-11}\,
\mathrm{m^3/(kg\,s^2)}\), por \(c_{\rm int}^3L^2/\hbar_{\rm ret}\).
& \(6.67430(15)\times10^{-11}\,\mathrm{m^3/(kg\,s^2)}\), CODATA 2022.\\
\bottomrule
\end{tabularx}
\caption{Evaluaciones y referencias en las realizaciones especificadas.
La fila térmica conserva la identificación de bases explicada en el texto.
La incertidumbre
relativa estándar de la referencia de \(G\) es aproximadamente
\(2.24743\times10^{-5}\); \(k_B\) carece de incertidumbre definicional.
La evaluación gravitatoria y su diferencia respecto de la referencia
utilizan la composición longitudinal y radial declarada en el texto.}
\label{tab:ii-kb-g-cartas}
\end{table}

Las fuentes externas de este último cuadro son las
\href{https://www.bipm.org/en/measurement-units/si-defining-constants}
{constantes definitorias del BIPM} y la
\href{https://physics.nist.gov/cuu/Constants/Table/allascii.txt}
{tabla CODATA 2022 del NIST} \cite{codata2022}.
La incertidumbre experimental de \(G\) no se transfiere a \(h\),
\(\hbar\) ni \(k_B\). Tampoco una igualdad de cifras en una coordenatización
elimina la necesidad de identificar el lector que produce la magnitud.

\subsection{Reproducción de las evaluaciones}

El programa \texttt{verificar\_comparaciones\_II.py}, incluido con las
fuentes, recibe las coordenadas internas almacenadas en archivos distintos
de los datos externos. Comprueba la composición del funcional de Barbero,
la conversión entre \(h\) y \(\hbar\), la raíz bibliográfica y su cola,
los módulos CKM, el cuarteto \(J\) y los residuos de esas comparaciones. Utiliza
aritmética decimal de 75 cifras y funciones elementales evaluadas después
de congelar las salidas internas. Las instrucciones de ejecución
y las versiones de los programas se incluyen en el paquete reproducible.
La precisión de trabajo no aumenta los dígitos disponibles del
cálculo de origen: se presentan únicamente cifras estables frente a su redondeo.
El control registra los archivos utilizados y sus huellas, de modo que
un cambio de evaluación o de edición externa quede visible.

La evaluación gravitatoria añadida se reproduce mediante
\texttt{verificar\_reproduccion.py}, en el directorio
\texttt{reproduccion\_gravitatoria\_20260926} de las fuentes.
El conjunto focal conserva las pruebas de transporte polar, composición
métrica, rigidez simultánea, acción de fase y composición conjunta,
junto con las sustituciones estructural y angular. Los verificadores
de sustitución leen las expansiones regionales archivadas y el registro
\(K\), reconstruyen las expresiones posteriores y comparan dos precisiones
de trabajo; las evaluaciones gravitatorias conservan 71 cifras
significativas estables. Las identidades generales se demuestran en el
cuerpo; los controles racionales examinan casos finitos de esas identidades.
Las entradas, la procedencia de las copias, sus huellas y el alcance de
cada programa están documentados en el mismo directorio. Esta ejecución
reutiliza las salidas generadas anteriores y conserva diferenciadas la
reproducción focal, la ejecución íntegra del TPK y la formalización Lean.

% NUMCHECK hbar_SI_coefficient=1.0545718176461563912 tol=1e-19
% NUMCHECK h_pre_coefficient=6.62607014999998793 tol=6e-18
% NUMCHECK h_ret_coefficient=6.62607014540625433 tol=6e-18
% NUMCHECK relative_pre=-1.822221e-15 tol=6e-22
% NUMCHECK relative_ret=-6.932836e-10 tol=6e-17
% NUMCHECK gamma_GM=0.2740668582328300 tol=6e-17
% NUMCHECK gamma_difference=-5.918553837e-5 tol=6e-15
% NUMCHECK gamma_relative=-2.159529202e-4 tol=6e-14
% NUMCHECK ckm.s12.value=0.225007404757 tol=6e-13
% NUMCHECK ckm.s23.value=0.041841757010 tol=6e-13
% NUMCHECK ckm.s13.value=0.003734601164 tol=6e-13
% NUMCHECK ckm.delta_rad.value=1.146265461116 tol=6e-13
% NUMCHECK ckm.s12.marginal_residual=-0.003817 tol=6e-7
% NUMCHECK ckm.s23.marginal_residual=0.014882 tol=6e-7
% NUMCHECK ckm.s13.marginal_residual=0.028902 tol=6e-7
% NUMCHECK ckm.delta_rad.marginal_residual=-0.028251 tol=6e-7
% NUMCHECK ckm.J.marginal_residual=-0.008564 tol=6e-7
% NUMCHECK gravity_selector=295.4525715010569 tol=6e-14
% NUMCHECK G_relative_reference_uncertainty=2.24743e-5 tol=6e-11
% END THERMAL_R03_SYNC_18
